\section{Discussion and Limitations}
\label{sec:discussion}

The updated results change the interpretation of OrbInspect from a preliminary
proxy-geometry demonstration to a mesh-aware high-coverage planning baseline.
The proposed set-cover CW tour reaches 98.33\% coverage with positive
clearance and the lowest $\Delta v$ among feasible methods that achieve the
same coverage. The comparison also clarifies why dynamics must be included in
view planning: the pure coverage-greedy method achieves the same coverage but
is infeasible because it violates the keep-out envelope. Coverage alone is
therefore not an adequate objective for orbital inspection.

The proposed method is not universally optimal across every metric. The
nearest-NBV baseline records a larger minimum clearance, and the proposed safe
CW NBV method uses fewer selected viewpoints. However, both require more
$\Delta v$ than the set-cover CW tour. This indicates that the current
weighted set-cover stage and CW tour ordering are effective for the chosen
mesh benchmark, but the weight settings should be studied across additional
geometries, initial states, target densities, and safety margins.

The mesh benchmark still uses a downsampled line-of-sight representation for
occlusion checking. This is appropriate for a reproducible paper baseline, but
a flight-oriented planner would need higher-fidelity geometry acceleration,
uncertainty-inflated keep-out volumes, camera slew limits, illumination
constraints, state-estimation error, plume constraints, and actuator
saturation studies. The reported RMS transfer tracking error is also a
planner-side rollout metric rather than a closed-loop flight validation error.

Gazebo validation should be interpreted conservatively. The videos and frames
verify that the visual model, chaser model, camera view, and pose-following
interfaces work in the integrated stack. They do not validate orbital
mechanics, prove coverage, or certify collision avoidance. In OrbInspect,
those quantitative claims come from the HCW dynamics, coverage, trajectory,
attitude, safety, and planner logs.

\subsection{Self-Contained Artifact Inventory}

All cited figures and data are contained in the paper folder. The key local
artifacts are listed in Table~\ref{tab:artifact-inventory}.

\begin{table}[t]
  \centering
  \caption{Self-contained manuscript artifacts.}
  \label{tab:artifact-inventory}
  \begin{tabular}{p{0.30\linewidth}p{0.56\linewidth}}
    \toprule
    Local artifact & Contents \\
    \midrule
    High-coverage figures &
    Mesh-overlaid trajectory, coverage, $\Delta v$, efficiency, and safety
    comparison figures. \\
    ISS mesh figures &
    High-resolution NASA ISS mesh preview in PDF, PNG, and SVG form. \\
    ISS mesh model &
    Copied NASA ISS GLB used to regenerate the mesh-overlaid trajectory
    figures. \\
    High-coverage data &
    Summary files, config snapshot, and raw comparison, planner, viewpoint,
    attitude, trajectory, and coverage CSV files. \\
    Gazebo figures &
    Gazebo preview frames and ROS smoke-run trajectory/control figures. \\
    ROS smoke data &
    Summary files, experiment manifest, config snapshots, and raw ROS
    execution CSV files. \\
    Legacy offline data &
    Earlier proxy and mesh single-run planner artifacts retained for
    traceability. \\
    Video data &
    Gazebo validation videos and preview frames. \\
    \bottomrule
  \end{tabular}
\end{table}
